\documentclass{article}
\usepackage{pathfinder-readable}

\title{Strategic Carbon Response: The Limit of a Direct Transfer}

\begin{document}
\maketitle

\begin{abstract}
Q studies payments that steer private resource requests towards a welfare optimum, while P forecasts
carbon intensity and schedules flexible electricity demand. The thread asked whether Q's mechanism could
make strategic participants carry out P's carbon-aware dispatch. The peers found a balanced charge for a
restricted, fixed-forecast problem and a simple test for whether a predicted emissions reduction survives
forecast error. Neither result covers P's coupled, partly binary dispatch, and a conflict between two of
Q's printed conditions further weakens a direct transfer. The verifier therefore paused the thread
pending a valid mechanism for P's constraints and counterfactual emissions evidence.
\end{abstract}

\section{Introduction}

Q considers people or organisations that request shares of limited resources and know their own
satisfaction better than the manager does. It proposes quadratic payments so that each participant's best
response leads to the allocation that maximises total satisfaction. Its stated results include a unique
Nash equilibrium, budget balance at that equilibrium, voluntary participation, and convergence of a
sample-based learning rule, each under stated assumptions. Its example concerns electric vehicle users
choosing routes and charging stations \cite{Q}.

P addresses a different problem in an electricity network. It forecasts the carbon intensity of power at
each node for the next day, then schedules flexible loads, including mobile energy storage systems and
distributed data centres, to reduce emissions. It tests the approach on a modified IEEE 33-bus network.
The ``agents'' in its forecasting workflow clean inputs and respond to forecast errors; they are not
modelled as self-interested owners of the scheduled loads \cite{P}.

The proposed connexion was to add private owners to P's dispatch and use Q's payments to make their
choices follow a carbon-aware plan. This raised two questions. Could P's decisions and constraints fit
Q's resource-allocation model? If they could, would a plan based on forecast carbon intensity still
reduce emissions after the changed dispatch altered power flows?

\section{What was done}

The peers first compared the two optimisation problems. Q allows continuous requests in convex local
sets, with a componentwise upper limit on their sum. P requires a workload balance across data centres
and uses binary decisions for storage locations and travel. The peers therefore set aside a direct
application to P's full model. They explored a smaller model in which owners allocate splittable workload
batches across destinations. That model can have convex local sets and destination capacities, but it
changes who the agents are and what they control \cite{Q,P}.

Next, one peer proposed a carbon charge for a fixed forecast; the other checked its algebra. They then
distinguished the forecast objective from emissions after dispatch, and derived an error test for
comparing a proposed dispatch with a baseline. Finally, they checked Q's printed payment and learning
rules and compared the proposed connexion with earlier work. The verifier's last check found an
additional conflict between Q's printed implementation and uniqueness conditions.

\section{Results within the restricted model}

Suppose there are \(N>1\) participants. Participant \(n\) chooses a nonnegative allocation vector
\(x_n\). Let \(\widehat e\) be a fixed vector of forecast carbon intensities and \(\beta\geq 0\) the
weight placed on forecast emissions. The peers checked the charge
\[
\tau_n^C(x)=\beta\widehat e^{\mathsf T}x_n
-\frac{\beta}{N-1}\sum_{m\ne n}\widehat e^{\mathsf T}x_m.
\]
Here \(m\) indexes the other participants. Each forecast term appears once as a charge and \(N-1\) times
in rebates, so the charges sum to zero for every allocation. A participant's own marginal charge is
\(\beta\widehat e\), because the rebate depends only on others' choices. Thus, \emph{if} a valid version
of Q's mechanism applies to this restricted allocation, adding the charge gives the same best responses
as replacing each private satisfaction function \(u_n(x_n)\) by \(u_n(x_n)-\beta\widehat e^{\mathsf
T}x_n\). The resulting objective trades satisfaction against forecast emissions; it is not P's pure
emissions objective. This is an algebraic construction, not an implementation result for P's dispatch.

The peers also obtained a comparison that does not depend on how a dispatch is chosen. Let \(C(x)\) be
emissions after power flows are recomputed for dispatch \(x\), and let \(\widehat C(x)\) be its forecast
value. For a proposed dispatch \(x^M\) and a baseline \(x^0\), suppose the absolute forecast error at
\emph{each} dispatch is at most \(\delta\). Then
\[
C(x^0)-C(x^M)
\geq \widehat C(x^0)-\widehat C(x^M)-2\delta.
\]
The argument is just the two error bounds: the baseline's realised emissions can be at most \(\delta\)
below its forecast, and the proposed dispatch's can be at most \(\delta\) above. A forecast reduction
greater than \(2\delta\) therefore certifies a realised reduction for those two decisions. P's historical
forecast error does not supply this error bound for dispatches chosen under a new incentive scheme. Nor
does maximising satisfaction less weighted forecast emissions, on its own, guarantee an emissions
reduction.

\section{Objections and unsettled points}

The first obstacle is the agent model. P's balance equality links data centres, while Q's shared
constraint is an upper capacity. Giving each workload batch its own choice over destinations could absorb
the balance into local constraints, but it creates a new market. A mandatory batch also rules out the
zero allocation used in Q's voluntary-participation argument. The nonnegative rebate in the charge above
is therefore insufficient to establish participation without a feasible outside option. P's binary
storage routes remain outside that restricted model \cite{Q,P}.

The second obstacle is the carbon signal. P's carbon-flow calculation makes nodal intensity depend on
dispatch, whereas its scheduling objective uses forecast intensities as fixed coefficients. A fixed
charge can match the forecast objective but need not match the marginal effect on realised emissions. The
peers' illustrative two-location example showed that a forecast can favour one location even when moving
all demand there makes realised emissions larger. This example was a logical check, not a fit to P's
network. P describes a closed-loop assessment and reports a 33.86\% reduction between its Cases 2 and 3;
the supplied text does not establish whether that result survives private costs and counterfactual
forecast error \cite{P}. The peers did not infer that P's reported comparison was wrong.

Checks of Q raised separate cautions. The reduced payment printed beside its learning update has an
own-price quadratic sign that differs from its full payment; a peer expanded both expressions and
recorded corrected blocks. Q's convergence theorem requires a growing sample size, but its numerical
schedule gives one sample at every iteration. These discrepancies call for correction and verification
before transferring its learning claim; they do not show that its plotted trajectories fail \cite{Q}.

The verifier identified a more direct conflict in Q's printed conditions. Let \(\Psi\) be Q's matrix of
payment curvature, let \(K\) be the number of resource coordinates, and let \(A_{nm}^{n}\) denote its
price-to-price block for participant \(n\) and participant \(m\). Q's implementation condition P2(i) says
\(\sum_m A_{nm}^{n}=0\). For any nonzero \(K\)-vector \(p\), take a direction \(v\) that changes every
participant's proposed price by \(p\) and leaves every resource request unchanged. The price blocks then
give \(v^{\mathsf T}(\Psi+\Psi^{\mathsf T})v=0\). Yet Q's uniqueness condition P1(ii) requires this
quantity to be at most \(-\upsilon\lVert v\rVert^2<0\), for a positive constant \(\upsilon\). As printed,
these two conditions cannot both hold. This does not rule out a corrected mechanism, but the thread
contains no such correction or replacement proof \cite{Q}.

The prior-work check also narrowed any novelty claim. Earlier work covers incentives for spatial
data-centre demand response \cite{chen_idc}, marginal emissions signals for demand response
\cite{ieee_marginal}, and carbon-aware scheduling under uncertain grid signals \cite{ieee_uncertain}. It
also covers equilibrium effects of average carbon signals \cite{jiang}, joint electricity-carbon pricing
with budget balance \cite{chen_zhao}, and carbon-flow-aware demand response \cite{chen_carbon_flow}. The
peers did not establish that the proposed combination is absent from this literature. The balanced charge
and the \(2\delta\) inequality are useful checks, but neither is a new result for P's full setting.

\section{Why the thread stopped}

The verifier paused the thread because the conditional charge and error test do not connect Q's stated
mechanism to P's coupled, partly binary dispatch; the prior-work check also left novelty unproved. The
conflict between Q's P2(i) and P1(ii) makes a direct appeal to its printed conditions untenable. The
smallest useful restart is a worked restricted data-centre case: specify the strategic owners and their
outside option, repair and verify the needed Q conditions or supply another mechanism for the balance
constraint, and measure forecast and realised emissions for its dispatch and a comparator on P's network.
A claim about the full storage and data-centre system would then still need to handle its routing
decisions.

\bibliographystyle{plainurl}
\bibliography{references}

\end{document}
